\chapter{Discussion et retour d'expérience}

A l'issue du développement et des tests, il convient de prendre du recul sur l'expérience menee. Ce chapitre proposé une lecture critique des deux paradigmes, identifié les forces et les limites de notre réalisation, et esquisse quelques pistes d'évolution.

\section{Comparaison effective des deux paradigmes}

\subsection{Verbosité des échanges}

Notre implémentation rend tangible la différence de verbosité entre les deux paradigmes. Une enveloppe SOAP \op{login} occupe environ trois cents octets, tandis que la requête REST équivalente, transmise au format JSON, ne depasse pas soixante octets. Sur des opérations fréquentes telles que la scrutation de messages, cet écart se traduit par une différence notable de bande passante consommee. A l'échelle d'un site grand public, cette economie peut justifier le choix de REST pour des raisons strictement pratiques.

\subsection{Contrat formel ou contrat implicite}

L'autre différence fondamentale tient à la nature du contrat. Le contrat WSDL d'un service SOAP est lisible par une machine : tout client peut interroger le WSDL et générer un stub fonctionnel sans intervention humaine. Notre service de chat Node.js consomme effectivement le WSDL du service Java par cette technique. À l'opposé, REST n'impose aucun contrat formel : la spécification est typiquement rédigée dans un document OpenAPI distinct, dont la cohérence avec l'implémentation réelle doit être vérifiée a part. Cette différence reflète des philosophies opposées : SOAP privilégie la rigueur, REST privilégie la souplesse.

\subsection{Outillage et intégration}

Du côté de l'outillage, REST profite de la disponibilite universelle des clients HTTP. Un simple appel \op{curl} suffit pour invoquer une opération, alors qu'un appel SOAP brut implique de rédiger soi-même l'enveloppe XML. Cette accessibilite explique en grande partie le succès de REST dans les écosystèmes ouverts. SOAP conserve toutefois un avantage dans les contextes d'intégration entre entreprises, ou la stabilité contractuelle est plus precieuse que la simplicité opérationnelle.

\subsection{Gestion des erreurs}

Les deux paradigmes traitent les erreurs differemment. SOAP impose une structure \texttt{Fault} dont la forme est standardisée : code, message, détail. REST réutilisé la semantique des codes de statut HTTP, complète par un objet d'erreur dans le corps de la réponse. Dans notre réalisation, ces deux mécanismes se sont reveles équivalemment expressifs ; aucun n'a manque pour traduire les exceptions métier.

\section{Forces de l'architecture réalisée}

\subsection{Démonstration concrète de la polyglossie}

Le choix d'implémenter les services en trois langages différents constitue, rétrospectivement, l'un des apports pédagogiques les plus marquants. La cohabitation harmonieuse de Java, Node.js et Python valide concrètement que les services Web tiennent leurs promesses d'interopérabilité : aucune des incompatibilités craintes n'est apparue, et chaque langage a su consommer naturellement les contrats produits par les autres.

\subsection{Hexagonalite vérifiable}

Le pattern hexagonal côté frontend a été éprouvé avec succès. Aucune des pages PHP applicatives ne contient de référence au protocole sous-jacent. La bascule entre SOAP et REST passe exclusivement par la fabrique \texttt{ClientFactory::make()}. Cette vérification empirique du pattern, possible dans un projet de taille modeste, illustre le bénéfice du découplage par abstraction.

\subsection{Composition de services observable}

L'isolation des bases de données a impose la composition de services. Cette contrainte, perçue initialement comme un obstacle, est devenue le coeur pédagogique du projet. Elle a permis de manipuler concrètement la notion de \emph{cohérence applicative} en lieu et place d'une cohérence relationnelle. Cette expérience prépare directement la transition vers une architecture de microservices.

\section{Limites de l'architecture réalisée}

\subsection{Polling et latence}

Le mécanisme de polling, simple à implémenter, présente une limite intrinsèque : la latence de livraison des messages est bornée par la période de scrutation, fixée à mille cinq cents millisecondes. Pour une chatroom pédagogique, ce délai est acceptable. Pour une application de messagerie professionnelle, il devrait être réduit, ce qui multiplierait la charge côté serveur et plaiderait pour une transition vers WebSocket ou Server-Sent Events.

\subsection{Absence de quotas et de limitation de débit}

Notre réalisation ne comporte aucun mécanisme de limitation de débit. Un utilisateur malveillant pourrait saturer le service de chat en envoyant des messages à haute fréquence. Dans un contexte de production, un middleware de type \emph{rate limiter}, place à l'entrée de chaque service, devrait être ajouté.

\subsection{Sécurité de transport}

Tous les échanges se font en HTTP non chiffre. Cette décision se justifie dans un contexte pédagogique de développement local, mais serait inacceptable en production. Un déploiement réel imposerait l'usage de TLS, éventuellement avec authentification mutuelle pour les appels inter-services.

\section{Polyglossie : retour d'expérience}

\subsection{Avantages observés}

La diversité des langages a force une discipline de conception : chaque service expose un contrat strict et réduit au strict nécessaire. Les conventions internes à chaque langage n'ont pas pu s'imposer aux autres, ce qui a abouti à des interfaces plus épurées que ce qu'aurait produit une réalisation monolangage.

\subsection{Surcouts mesurés}

La polyglossie introduit toutefois des surcoûts. La mise en place de l'environnement de développement multiplié les outils à maîtriser : Maven pour Java, npm pour Node, pip pour Python, Composer ou serveur intégré pour PHP. Le diagnostic d'une erreur traversant plusieurs services impose de consulter plusieurs journaux applicatifs au lieu d'un seul. Ces surcoûts, acceptables dans un projet pédagogique, doivent être évalués lucidement avant d'être transposés dans un contexte industriel.

\section{Perspectives}

\subsection{Vers une notification temps réel}

La première évolution naturelle consisterait à remplacer le polling par un canal poussé. WebSocket, standardisé par la RFC 6455, conviendrait particulièrement aux échanges bidirectionnels d'une chatroom. Cette évolution toucherait principalement le frontend et un nouvel endpoint côté chaque \service{ChatService}, sans remettre en cause le reste de l'architecture.

\subsection{Adoption de gRPC pour les communications inter-services}

Pour les appels internes entre services, gRPC offrirait une alternative intéressante a SOAP et REST. Fondé sur HTTP/2 et sur la serialisation Protocol Buffers, il combine la rigueur contractuelle de SOAP, la performance de REST, et un support natif du streaming bidirectionnel. Une troisième phase \emph{gRPC} de notre application constituerait un prolongement cohérent du présent travail.

\subsection{Conteneurisation complète des services applicatifs}

Notre infrastructure actuelle ne conteneurise que PostgreSQL. Une évolution vers une conteneurisation complète, éventuellement orchestree par Docker Compose ou Kubernetes, simplifierait le déploiement et rapprocherait l'environnement de développement d'un environnement de production réaliste.

\subsection{Documentation API automatisee}

La phase REST de notre application beneficierait grandement d'une spécification OpenAPI générée directement depuis le code. Du côté Java, des bibliothèques telles que swagger-jaxrs2 produisent cette spécification automatiquement. Du côté Python, Flask-Smorest ou Flask-RESTX offrent une fonctionnalite équivalente. Cette documentation autogeneree faciliterait l'intégration de futurs clients.

\section{Synthèse}

Notre expérience confirme que les deux paradigmes coexistent legitimement dans l'écosystème actuel des services Web. REST l'emporte sur la simplicité et la souplesse, SOAP conserve l'avantage de la rigueur contractuelle. La connaissance des deux est un atout pour tout ingénieur en charge d'integrations distribuées, qu'il s'agisse de concevoir un nouveau service ou de dialoguer avec un système legacy. Le chapitre suivant tiré les enseignements généraux de ce projet.
